Let the integer have $d$ digits. Since its reversal has the same length, multiplying by $9$ forces the leading digit of the original to be $1$. The last digit must be $9$, since $9b\equiv1\pmod{10}$. A one-digit number cannot work.
The two-digit candidate $19$ fails. A three-digit candidate $1a9$ would satisfy $9(109+10a)=901+10a$, giving $a=-1$, impossible.
For four digits write $1ab9$. The equation
\[
9(1009+100a+10b)=9001+100b+10a
\]
reduces to $b=89a+8$. Digit bounds force $a=0,b=8$. Thus $\boxed{1089}$ works, with $9\cdot1089=9801$, and no smaller number does.
